\chapter{Technical appendices}
\label{ch:expanded-appendices}

The main chapters make the investment case in the order a decision-maker
needs: an observed problem, a mechanism, a bounded product, evidence that
supports parts of the mechanism, and a test that can change the decision.
This appendix supplies the working detail behind that route. It is written
for the engineer who will build a fixture, the researcher who will audit a
claim, and the sponsor who wants to know exactly what a successful twelve
weeks would leave behind.

The appendices are deliberately explanatory. A table records a contract, but
the paragraphs around it explain why the contract exists, what a failure looks
like, and what conclusion the failure permits. The examples use LaTeX because
it is the richest current target in the repository; the same objects can be
represented in Markdown or a word-processing format once their source and
rendered anchors are known.

\section{Notation and terms}
\label{app:notation}

\subsection{The unit of work}

Humanvoice operates on a \emph{document run}. A run is one immutable source
snapshot, one declared reading brief, one execution configuration, and one
rendered result. A new commit, a changed bibliography, a different parser, or
a different model version creates a new run. This simple rule prevents an
author from comparing a revised page with a stale source map and then calling
the result a protected diff.

Let a run be
\[
  R=(S,V,C,E,P),
\]
where $S$ is the source snapshot, $V$ its version identifier, $C$ the reading
brief, $E$ the execution environment, and $P$ the rendered pages and
anchors. The tuple is not a claim that every field must be stored forever. It
is the minimum information needed to reproduce a finding while the decision
is still live. A retention policy may later delete page images while keeping a
hash and a decision record; it must then say that it has done so.

The reading brief has eight fields. The first four define the reader's task;
the remaining four keep later inspections from quietly changing that task.

\begin{enumerate}
\item the reader role, such as investment committee member, technical
      referee, or policy executive;
\item the decision the reader must be able to make after reading;
\item the knowledge and time budget that may reasonably be assumed; and
\item the evidence, uncertainty, and disclosure obligations that constrain
      the decision.
\item two or three named exemplar documents that establish the desired pace
      and explanatory register;
\item vocabulary the reader can reasonably be expected to know without a
      local gloss;
\item declared exemption zones, such as a map, decision box, or reference
      table, where a name may appear before it is taught; and
\item the versioned reader directives and the date of the last reader
      session.
\end{enumerate}

The brief is not a marketing persona. It is a testable statement about
what the document owes a named class of reader. For example, ``an expert
reader'' is too broad for an acceptance endpoint. ``A former econometrics
professor deciding whether to fund a twelve-week feasibility build,
with twenty minutes for the executive route and access to the appendices'' is
specific enough to generate a rubric and a timed exercise.

The exemplar field is deliberately comparative. It does not turn one
monograph into a universal standard. A pacing meter reports where the target
falls relative to the supplied exemplars and shows the passages that account
for the difference. The vocabulary and exemption fields prevent a first-use
audit from treating a reader's known language or an announced map as a defect.
The brief is cumulative: a reader's stopping point or standing instruction
updates the next version, with the change recorded rather than inferred from a
new score.

\subsection{Pacing, ordering, and assertion records}
\label{app:pacing-records}

The HPO case suggests three related but non-interchangeable records. A
\emph{pacing profile} counts concepts and paragraphs; a \emph{first-use
record} checks whether a term was introduced before it was needed; an
\emph{assertion audit} checks whether an editorial repair smuggled in a new
claim. All three locate questions. None can certify that a reader will accept
the page.

For a document unit $u$, let $W_u$ be prose words and $C_u$ the distinct
concepts first introduced in that unit. The descriptive introduction rate is

\begin{equation}
  q_u=\frac{W_u}{C_u},
  \qquad
  b_u=\frac{\#\{\text{paragraphs introducing at least three concepts}\}}
                 {\#\{\text{prose paragraphs}\}}.
  \label{eq:pacing-profile}
\end{equation}

The implementation also reports singleton share, repeated occurrences, and
the number of first introductions in a trailing 1,500-word window. These
quantities are diagnostic: the same detector must run over the target and its
exemplars, and the author must see the spans behind any comparison.

\begin{table}[H]
\centering
\small
\begin{tabularx}{\textwidth}{@{}p{0.22\textwidth}L@{}}
\toprule
Record & Minimum fields \\
\midrule
Pacing profile & run id; unit; detector version; prose-word count; concept
count; words per concept; singleton share; burst paragraphs; live-load
window; exemplar ids; relative position; dismissed findings. \\
First-use record & concept; first source line and rendered page; introduction
signal; known-vocabulary match; exemption-zone match; taught-later pointer;
author disposition. \\
Assertion audit & parent and child run; added sentence span; detected number,
comparison, superlative, or causal verb; citation or protected-claim pointer;
author explanation; status. \\
Reader observation & reader role; stopping page or section; quoted span;
missing relation; elapsed time; requested change; confidence. \\
\bottomrule
\end{tabularx}
\caption{Small records that turn the HPO lessons into inspectable questions.}
\label{tab:pacing-order-assertion-records}
\end{table}

\subsection{Objects, spans, and anchors}

The source contains more than prose. Humanvoice calls every meaning-bearing
or navigation-bearing item an \emph{object}. A span is a contiguous source
range; an anchor is a stable identifier that connects a source range to a
rendered location. The same object may have several spans: a citation key in
the sentence, its bibliography entry, and the rendered citation are one
identity relation rather than three independent strings.

For an object $o$, write
\[
  o=(\operatorname{kind},\operatorname{span},\operatorname{anchor},
     \operatorname{owner},\operatorname{status}).
\]
The owner is the person or rule allowed to change it. A numerical literal may
be owned by the analyst; a quotation may be owned by its source; a claim
scope may require the document author; a page anchor is owned by the build
system. Ownership is useful because it turns a vague warning into a bounded
decision: the tool can flag a change, but it cannot silently decide that the
owner's object no longer matters.

The first implementation should distinguish these object kinds:

\begin{table}[H]
\centering
\small
\begin{tabularx}{\textwidth}{@{}p{0.19\textwidth}p{0.23\textwidth}L@{}}
\toprule
Kind & Typical anchor & Why it matters \\
\midrule
Prose claim & source line and paragraph id & A polished sentence can widen a
claim even when every word is grammatical. \\
Assumption & tagged span or equation premise & A reader must know which
conditions make the conclusion conditional. \\
Equation & equation label and source span & Formatting changes can hide a
changed term or a lost restriction. \\
Number and unit & token range & A decimal or unit change can alter the
decision while escaping a prose diff. \\
Citation identity & citation key plus bibliography id & A citation that still
looks plausible may resolve to a different source. \\
Cross-reference & label and target page & Navigation is part of a long
document's argument. \\
Quotation or code & fenced or delimited span & Exactness and provenance are
different from ordinary copy-editing. \\
Table or figure & environment id and caption & A result can be correct in the
text but misaligned with its displayed evidence. \\
\bottomrule
\end{tabularx}
\caption{Initial meaning-bearing object vocabulary.}
\label{tab:object-vocabulary}
\end{table}

The object vocabulary is intentionally smaller than a complete abstract
syntax tree. A parser may know many more nodes, but the product should expose
only those that affect a reader decision or a protected comparison. This is a
product boundary, not a limitation on internal implementation. It keeps the
author's review queue legible and makes test coverage auditable.

\subsection{Finding, decision, and revision}

A \emph{finding} is a located, falsifiable prompt about an object or relation.
It is not a replacement sentence. Formally, a finding can be represented as
\[
  f=(R,\operatorname{kind},\operatorname{location},\operatorname{evidence},
     \operatorname{question},\operatorname{confidence},\operatorname{state}).
\]
The evidence field points to a source span, a deterministic count, a reader
observation, or a declared absence. The question is phrased so that an author
can answer ``keep'', ``change'', or ``abstain'' without accepting an implied
rewrite. Confidence describes the diagnostic's reliability, not the truth of
the underlying claim.

The author decision is a separate record:
\[
  d=(f,\operatorname{action},\operatorname{rationale},\operatorname{new\ run}).
\]
An accepted finding may lead to a revision; a rejected finding remains useful
because it teaches the diagnostic about this author and genre. An abstention
means the product could not safely decide. Counting rejection and abstention
is essential: a tool that raises fewer flags by hiding uncertainty can appear
precise while making the author less safe.

\begin{cardexample}{A small run, written out}
The source sentence says, ``The intervention improves quality.'' The reader
contract names a referee who must decide whether to fund a feasibility test.
The evidence record shows two internal repair cases, neither with a controlled
comparison. Humanvoice therefore emits a finding at the word ``improves'':
``Which observed outcome supports this comparative verb, and should it be
narrowed to a testable hypothesis?'' The author changes the sentence to
``The intervention is designed to test whether review burden falls.'' The
equation, citation identities, and decision request remain unchanged. The
finding is closed with the rationale ``claim changed from result to hypothesis''.
\end{cardexample}

The tempting but invalid inference is that the tool has proved the revised
sentence better. It has only made the claim class explicit and preserved the
objects that must be reviewed. A reader study is still required for an
effectiveness claim.

\subsection{Claim classes and evidence grades}

Every load-bearing sentence in the proposal should be classifiable as one of
four kinds: source result, local implication, project hypothesis, or open
question. The classification is a reading aid, not a substitute for a
bibliography. A source result needs a citation and an inspectable anchor. A
local implication must state the bridge from the source result to the product.
A project hypothesis must name the planned observation. An open question must
name what would resolve it.

For implementation decisions, use an evidence grade with two axes. The first
axis is \emph{provenance}: primary external study, official package record,
internal case, or design assumption. The second is \emph{fit}: direct fit to
the target population, adjacent fit, or unknown fit. A highly credible study
with adjacent fit should not outrank a modest internal observation when the
question is simply whether a source map survives this repository's macros;
neither should be used to support an efficacy claim outside its scope.

\begin{table}[H]
\centering
\small
\begin{tabularx}{\textwidth}{@{}p{0.19\textwidth}p{0.28\textwidth}L@{}}
\toprule
Evidence grade & Interpretation & Permitted use \\
\midrule
P--D & Primary or official source, direct target fit & Support a bounded
design or benchmark expectation. \\
P--A & Primary or official source, adjacent fit & Motivate a test; do not
generalize the outcome. \\
I--D & Internal record, direct repository fit & Specify a fixture or
workflow requirement; not external efficacy. \\
I--A & Internal record, adjacent or selected case & Generate a hypothesis or
counterexample. \\
A--U & Assumption with unknown fit & Price or scope the experiment only. \\
\bottomrule
\end{tabularx}
\caption{A compact evidence vocabulary for design decisions.}
\label{tab:evidence-grades}
\end{table}

\section{Detailed requirements and records}
\label{app:requirements}

\subsection{From a promise to an acceptance contract}

The product proposal uses requirements in a narrow sense: a requirement is a
behavior that can be observed in a run and accepted or rejected by a named
owner. ``The writing should feel human'' is a valuable aspiration but not a
first-release requirement. ``A reviewer can open a finding and reach the
source span and rendered page within two interactions'' is testable. The
distinction prevents a broad editorial goal from being hidden behind a
surface-style score.

This is the canonical R1--R24 acceptance register. The implementation contract
is in Appendix~\ref{app:implementation-contract}, with machine-readable
schemas in \path{schemas/}. The contract annex defines execution, schemas,
exit codes, and authority without creating a second release register.

\begin{longtable}{@{}p{0.07\textwidth}p{0.19\textwidth}p{0.40\textwidth}p{0.20\textwidth}@{}}
\caption{Canonical implementation requirements and acceptance observations.}
\label{tab:expanded-requirements}\\
\toprule
ID & Requirement & Acceptance observation & Owner \\
\midrule
\endfirsthead
\toprule
ID & Requirement & Acceptance observation & Owner \\
\midrule
\endhead
R1 & Reader-facing separation & A low-context reader can state purpose, object,
and requested action; internal register language is absent; the released file
traces to a source hash. & Domain editor \\
R2 & Policy precedence & Conflict fixtures resolve correctness, source fidelity,
meaning, comprehension, then regularity; format overrides are configuration. &
Product engineer \\
R3 & Located diagnostics & Fixtures return stable spans and no inspection
command mutates input; accepted edits are separate records. & Product engineer \\
R4 & Layered diagnosis & Replay retains register, repetition, defensive,
substance, and comprehension results in order. & Evaluation lead \\
R5 & Protected diff & Equations, labels, citations, numbers, claims, code,
tables, and quotations have no unexplained material change. & Product engineer \\
R6 & Citation support & Identity and claim support are separate; unresolved
load-bearing support blocks release. & Domain editor \\
R7 & Calibrated judging & Frozen human labels report agreement, order,
verbosity, and family sensitivity before model judging is used. & Evaluation lead \\
R8 & Human acceptance & A promoted version records the pseudonymous reader,
hash, rubric, decision, requested changes, and elapsed time. & Independent reader \\
R9 & Held-out evaluation & Tuning and holdout hashes are frozen and no smoke
test is reported as effectiveness. & Evaluation lead \\
R10 & Minimal process surface & Provenance remains in the packet record while
the manuscript contains no orchestration state. & Domain editor \\
R11 & Privacy and disclosure & Unapproved remote routes are blocked; approved
calls record provider, model, purpose, retention, and approval. & Security or
policy owner \\
R12 & Aggregate convergence & Drift analysis refuses single-document
authorship classification, suppresses small cells, and reports uncertainty. &
Evaluation lead \\
R13 & Interaction telemetry & Intent, suggestion, acceptance, edit, rejection,
and rollback replay without unnecessary manuscript retention. & Product engineer \\
R14 & Heterogeneous outcomes & Prespecified writer, genre, and reader strata
report uncertainty and apply the harm stop rule. & Evaluation lead \\
R15 & Parser contract & Candidate adapters emit the same source-map and
protected-span contract and pass the locked malformed-input fixtures. & Product
engineer \\
R16 & Package adoption & Availability, operational fit, and effectiveness are
separate; adoption requires a held-out benchmark and independent review. &
Evaluation lead \\
R17 & Complete authoring brief & Missing critical reader, decision, evidence,
privacy, or protection fields fail before writer invocation; a complete brief
has a stable hash. & Document owner \\
R18 & Argument blueprint & Missing warrants, early concepts, duplicate jobs, or
purposeless figures block promotion; the plan has a stable hash. & Domain editor \\
R19 & Bounded drafting & Each unit records approved context and evidence and
cannot advance while its checks are unresolved. & Product engineer \\
R20 & Fail-closed release & Every seeded gate failure emits no reader packet and
preserves a located repair record. & Product engineer \\
R21 & Construct-specific critics & Held-out false-positive, false-negative,
order, verbosity, and family results are recorded before release use. &
Evaluation lead \\
R22 & Mutation testing & Critical deterministic mutations are caught or
abstained from; rhetorical mutation uncertainty is reported. & Evaluation lead \\
R23 & Bounded repair & Named strategies converge or stop with an unresolved
author choice; no loop reaches the reader silently. & Product engineer \\
R24 & Attention budget & Feasibility reports release rate, author and reader
minutes, abstentions, and meaning errors against the incumbent. & Sponsor or
project owner \\
\bottomrule
\end{longtable}

The register is the single release view. The machine-readable audit protocol
adds evidence links and implementation phases without changing these
behaviors. A result is reported by requirement family and phase; it is never
averaged into one release score. R5, R11, R17, R18, and R20 are hard blockers
for the protected core or full path respectively, as specified in the
implementation contract.

\subsection{Record schemas}

The record model should remain plain enough to inspect with ordinary command-
line tools. A JSON implementation is sensible, but the semantics matter more
than the serialization. The following fields are sufficient for the MVP.

\begin{table}[H]
\centering
\small
\begin{tabularx}{\textwidth}{@{}p{0.21\textwidth}L@{}}
\toprule
Record & Required fields and interpretation \\
\midrule
Run & run id; source hash; parent run; source format; build command; tool and
model versions; privacy class; created and completed times; outcome. \\
Object & object id; kind; source span; normalized value; rendered anchor;
owner; protection rule; comparison status. \\
Finding & finding id; run id; category; location; evidence; question;
confidence; abstention reason; state; diagnostic version; unit; repair move;
dismissal reason. \\
Decision & finding id; action; author id or role; rationale; parent and child
run ids; time spent; reviewer comments. \\
Reader response & document run; reader role; contract version; completion
time; comprehension answers; credibility rating; stopping point; quoted span;
missing relation; requested action; confidence; disclosure acknowledgement. \\
Pacing profile & unit; detector version; target metrics; exemplar metrics;
relative position; burst locations; dismissed concepts. \\
First-use audit & concept; first-use location; introduction signal; whitelist
or exemption match; taught-later pointer; disposition. \\
Assertion audit & parent and child run; added span; assertion type; support
pointer; author disposition. \\
Benchmark result & fixture id; candidate; version; expected object set;
observed object set; precision; abstention; runtime; failure trace. \\
\bottomrule
\end{tabularx}
\caption{MVP records and the questions they answer.}
\label{tab:record-schemas}
\end{table}

The schema separates identity from content. A source hash identifies what was
run; a rendered PDF is an output of that run, not its identity. This matters
when a build tool changes page breaks without changing source semantics. The
comparison can then say ``same source objects, different pagination'' rather
than reporting a false substantive drift.

\subsection{Traceability without a bureaucracy}

Traceability is useful when it follows the reader's question. For each
requirement, link one fixture, one observed result, and one decision rule. Do
not create a separate narrative for every administrative event. A compact
trace looks like
\[
  \text{requirement}\;\longrightarrow\;\text{fixture}\;\longrightarrow\;
  \text{result}\;\longrightarrow\;\text{release decision}.
\]
For R5, the fixture contains an equation whose denominator changes in a
revision; the result reports whether the change was surfaced; the decision
rule blocks a protected comparison if the denominator is unparsed. For R8, the
fixture is a timed reader exercise; the result is a response and a completion
time; the decision rule says whether the reader endpoint is observable enough
for the comparative study.

\begin{checkpoint}{The traceability test}
Ask an engineer to pick any promised behavior and an executive to pick any
investment conclusion. The engineer should be able to walk forward to a
fixture and a release rule. The executive should be able to walk backward to
the observed evidence and its limitations. If either path ends at a slogan,
the proposal is not yet operational.
\end{checkpoint}

\section{Internal case evidence and repair-pair index}
\label{app:case-index}

The repository contains a rare kind of product evidence: paired passages in
which a writer or editor changed a document after a reader failure. These
pairs are not randomized experiments. They are valuable because they preserve
the causal story that a feature must eventually support: what the reader
could not do, what changed, and what remained protected.

\subsection{Case selection and provenance}

The cases came from work on CardNPV, BGS, SMEwallet, BayesFilter with the
zero-lower-bound rescue, and MacroFinance. Each has a different scale and
failure mode. CardNPV supplies the strongest teaching example because its
small worked example and later abstraction make the reader's reconstruction task
visible. BGS supplies a negative process case: a long sequence of governance
and review activity did not produce an accepted document. SMEwallet supplies
a positive local gate. BayesFilter supplies a layered technical argument in
which notation and computational commitments must survive exposition.
MacroFinance supplies the book-scale problem: repetition and infrastructure
can be useful to an implementer while becoming a burden to a first reader.

The case index should preserve four provenance fields: source path or commit,
date of the observed reader interaction, the role of the person who made the
judgment, and whether the passage was later used in another revision. A
passage copied into several drafts is one lineage, not several independent
observations.

\begin{longtable}{@{}>{\raggedright\arraybackslash}p{0.15\textwidth}
>{\raggedright\arraybackslash}p{0.19\textwidth}
>{\raggedright\arraybackslash}p{0.25\textwidth}
>{\raggedright\arraybackslash}p{0.25\textwidth}@{}}
\caption{Internal case index and the claim boundary for each case.}
\label{tab:case-index}\\
\toprule
Case & Reader task & Failure or repair pair & What it can establish \\
\midrule
\endfirsthead
\toprule
Case & Reader task & Failure or repair pair & What it can establish \\
\midrule
\endhead
CardNPV & Reconstruct a finite value question and decide what is identified. &
Finite teaching case, ledger, timing, mechanism, and later abstraction. & A
training route can make objects and transitions inspectable; not a causal
estimate of Humanvoice. \\
BGS & Decide whether a proposal has a defensible purpose and next action. &
Repeated governance prose versus located purpose, evidence, and decision. & A
reader-work failure can persist despite compilation and multiple reviews. \\
SMEwallet & Pass a named reader gate for a business-facing technical unit. &
Local explanation and revision until the reader can restate the action. & A
unit-level acceptance record is feasible; not a population effect. \\
BayesFilter--ZLB & Follow assumptions through filtering and a boundary-case
rescue. & Layered notation, an uncovered failure, and a repair that preserves
the mathematical object. & Protected technical exposition needs more than
sentence-level style. \\
Macro\\Finance & Navigate a long methods and implementation volume. & Repeated
structure and source infrastructure made the book usable to builders but
harder to enter as a reader. & Monograph-scale navigation and appendices are
product requirements. \\
\bottomrule
\end{longtable}

\subsection{Repair-pair taxonomy}

A repair pair should be annotated by the first reader action it enables. The
following taxonomy is more informative than labels such as ``awkward'' or
``AI-like''.

\begin{enumerate}
\item \textbf{Orientation repair:} the opening identifies the problem,
      decision, and route before terminology accumulates.
\item \textbf{Actor repair:} an abstract process is assigned to a concrete
      author, reader, institution, or data-generating mechanism.
\item \textbf{Evidence repair:} a conclusion is connected to an observation,
      citation, calculation, or explicit absence of evidence.
\item \textbf{Scope repair:} a result is narrowed to the population, period,
      document unit, or claim class actually supported.
\item \textbf{Sequence repair:} definitions and examples are reordered so
      that each later step uses an object the reader has already seen.
\item \textbf{Register repair:} administrative or defensive language is
      replaced by the domain statement it was concealing.
\item \textbf{Navigation repair:} labels, headings, tables, and transitions
      make a long argument recoverable after interruption.
\item \textbf{Disclosure repair:} assistance, uncertainty, and ownership are
      stated where a reader's interpretation can be affected.
\end{enumerate}

For each pair, record the smallest change that enabled the action and the
objects deliberately left unchanged. This last field matters. A passage that
reads better only because its equation, qualification, or numerical result
was silently removed is not a safe repair.

\begin{table}[H]
\centering
\small
\begin{tabularx}{\textwidth}{@{}p{0.19\textwidth}p{0.28\textwidth}L@{}}
\toprule
Annotation & Positive observation & Counterexample to retain \\
\midrule
Orientation & A reader unfamiliar with the project states the decision after the opening. &
The opening is elegant but the reader still cannot name the object being
decided. \\
Evidence & The reader can point to the calculation or source supporting the
claim. & A citation is present but its result does not entail the sentence. \\
Scope & The revised claim names its population and uncertainty. & The prose is
qualified so heavily that no testable claim remains. \\
Sequence & A reader predicts why the next section is needed. & A glossary
defines every term but no transition explains their relationship. \\
Protection & The equation and number survive the prose repair. & A textual
diff is clean because the mathematical object was dropped. \\
\bottomrule
\end{tabularx}
\caption{Repair annotations should retain negative examples.}
\label{tab:repair-annotations}
\end{table}

\subsection{What the cases do not show}

The cases are selected for learning, not prevalence. They do not show how
often expert documents fail, whether a particular diagnostic causes the
repair, or whether the same repair works for a different reader. They also
contain author and editor judgment that cannot be reconstructed perfectly from
the final PDF. The feasibility corpus should therefore include both these
historical pairs and prospective cases collected under a fixed protocol.

The tempting but invalid inference is ``five successful repairs imply a
reliable product.'' The defensible conclusion is narrower: the repository
contains enough varied failures to specify fixtures, reader tasks, and
protected objects. That is precisely the kind of evidence a product proposal
needs before it asks for an efficacy claim.

\section{The HPO density case and its fixture}
\label{app:hpo-density-record}

The HPO survey engagement is the clearest internal example of why the product
needs more than a surface linter. One manager reviewed a 90-page methods
monograph over three rounds. The manager used the same word, ``dense,'' in the
first two rounds and then supplied a more precise ordering instruction in the
third. The agent's self-review said the document was readable before every
round; the manager's located observations disagreed.

The case record is dated 24--25 August 2026. Its numbers come from the
engagement ledger and the prototype concept-density script in the HPO survey
repository. They are internal observations with direct fit to the proposed
workflow, not external literature and not a representative sample. We retain
them because they specify fixtures and questions that a later independent
review can replay.

\begin{longtable}{@{}p{0.16\textwidth}p{0.25\textwidth}p{0.26\textwidth}p{0.22\textwidth}@{}}
\caption{HPO case record and evidence boundary.}
\label{tab:hpo-case-record}\\
\toprule
Round & Observed complaint & Located evidence & Permitted use \\
\midrule
\endfirsthead
\toprule
Round & Complaint & Evidence & Permitted use \\
\midrule
\endhead
1 & ``Too dense.'' & One 210-word paragraph contained six mechanisms and five
citations; mean prose-paragraph length moved from 128 to 90 words after repair.
& Seed paragraph-packing fixture; not a paragraph-length standard. \\
2 & ``Too dense.'' & Words per new concept were 70 in the survey and 83 and 253
in the two supplied exemplars; paragraphs introducing at least three new
concepts were 12.9, 5.9, and 0.8 per cent. & Calibration comparison; not a
universal threshold or efficacy result. \\
3 & ``Do not use it before explaining it.'' & Approximately 25 adjudicated
first-use violations, including an unexplained running-example algorithm and
early acquisition-function names. & Ordering fixture; prevalence outside this
case is unknown. \\
\bottomrule
\end{longtable}

The case also supplies a useful negative set: 80 passages that the first-use
audit flagged but the manager adjudicated as false alarms. The proposed
fixture stores the 25 positive locations and the 80 negative locations without
mixing them into rule tuning. Each item carries the source line, the first-use
term, the reader's known-vocabulary decision, any exemption-zone declaration,
the adjudicated label, and the repair or dismissal reason. A future precision
estimate must be made on a partition that the rule author has not inspected.

The concept detector is a standard-library prototype. It recognises a small
set of signals---acronyms, mid-sentence proper terms, emphasized definition
phrases, recurring technical bigrams, and citation keys---and removes common
LaTeX environments before counting. It has no sentence-level dependency
parser. The HPO engagement spot-checked roughly 80 percent precision with
noise that appeared similarly across the target and exemplars; Humanvoice
should replicate that estimate before treating the detector as an operational
component. Dismissal must be cheap, and a noisy list must never block an
author's ordinary edit.

The fixture therefore has four roles:

\begin{enumerate}
\item a \textbf{calibration role} for choosing wording, displays, and
      exemption conventions;
\item a \textbf{held-out role} for measuring precision and abstention after
      those choices are frozen;
\item a \textbf{reader-study role} for recording stopping points and quoted
      spans rather than relying on a score; and
\item a \textbf{claim-boundary role} that prevents one successful engagement
      from being described as a general reduction in review rounds.
\end{enumerate}

The implementation loop is consequently small and repeatable: select one
chapter, run the meter and first-use audit, inspect the ranked passages, let a
human author repair or dismiss them, rebuild the chapter, and record the
reader's stopping point. Only after that unit is stable should the same rule
be propagated to the next chapter. The ledger records the unit and brief
version so a later reader can tell which instruction was active at the time.

\section{Literature claim-to-design matrix}
\label{app:literature-matrix}

The literature chapter explains the main results in prose. This appendix gives
the audit-friendly matrix that keeps a source result distinct from the design
choice that follows it. The rows are grouped by question, not by publisher or
model family, because the product decision is comparative.

\subsection{How to use the matrix}

For each row, the researcher records the source identity, the inspected text
or official documentation, the population and task, the result, its limit,
and the Humanvoice implication. A source can appear in more than one row when
it answers different questions, but each implication must remain local. The
matrix is not a ranking of papers. It is a guard against a common citation
error: using evidence about model output style as if it were evidence about
reader acceptance.

Figure~\ref{fig:traceability-chain} summarizes the relationship that the rows
below must preserve. A source earns influence only through an inspectable
claim, design choice, and test.

\hvFigureTraceability

\begin{longtable}{@{}>{\raggedright\arraybackslash}p{0.14\textwidth}
>{\raggedright\arraybackslash}p{0.22\textwidth}
>{\raggedright\arraybackslash}p{0.25\textwidth}
>{\raggedright\arraybackslash}p{0.24\textwidth}@{}}
\caption{Claim-to-design matrix for the evidence review.}
\label{tab:literature-claim-design}\\
\toprule
Question & Source result to verify & Product implication & Open test \\
\midrule
\endfirsthead
\toprule
Question & Source result to verify & Product implication & Open test \\
\midrule
\endhead
Why outputs look alike & Distributional markers, repeated discourse, or
model-conditioned style are observed in a stated corpus. & Use patterns as
diagnostic prompts and compare against a genre baseline. & Do the prompts
reduce reader work without erasing legitimate disciplinary style? \\
Can detection identify origin? & Detection performance varies with domain,
editing, and generator; false positives remain consequential. & Never expose
a per-document origin verdict as the product outcome. & Does disclosure plus
located evidence improve a reader's decision? \\
Does feedback help? & Controlled writing or editing studies report changes in
time, quality, or task completion under specific conditions. & Test bounded
feedback with a declared incumbent and reader endpoint. & Does the effect
survive expert technical documents and longer tasks? \\
How do people use models? & Users accept, reject, and adapt suggestions
heterogeneously; overreliance and verification costs matter. & Record author
decisions, burden, abstention, and protected-object incidents. & Which reader
and writer populations benefit or are harmed? \\
Can models judge text? & Agreement can be high on surface edits and lower on
scope, evidence, and technical meaning; order and verbosity biases occur. &
Use model judges only as calibrated screening instruments with human labels. &
Does calibration hold on held-out mathematical and policy cases? \\
What is readability? & Formula scores depend on genre, audience, and text
features; comprehension is not identical to ease. & Pair any metric with a
reading brief and timed reconstruction task. & Which measures predict the
decision the document is meant to support? \\
What are integrity risks? & Provenance, disclosure, privacy, and authorship
norms differ by venue and can change behavior. & Make processing boundary and
assistance visible; preserve source ownership. & What policy applies to each
target institution and publication? \\
\bottomrule
\end{longtable}

\subsection{Search coverage and omitted-source risk}

A credible survey reports not only what it found but also where it may be
thin. The current review covers the supplied internal material, a targeted
set of primary studies on generative assistance and feedback, readability and
technical communication, model-based evaluation, and official documentation
for candidate software. It does not yet constitute a registered systematic
review. Specialist work in human--computer interaction, writing studies,
information integrity, accessibility, and privacy may change the design.

The next evidence pass should use at least two independent screeners, a
recorded query set, backward and forward citation chasing for load-bearing
studies, and a disposition for every retrieved item. A source should be
included because it answers a declared question, not because its conclusion
supports the investment. The screening sheet should preserve exclusions and
their reasons; otherwise the final bibliography will look comprehensive while
the selection process remains invisible.

The register below names the present risks one by one. ``Candidate'' means that
the source has been identified but has not been used to support a manuscript
claim. ``Included, bounded'' means that the manuscript states the source's
limited contribution while appraisal or independent verification remains
open. This is the omission register promised in the main text, not a list of
topics that someone might search later.

\begin{longtable}{@{}>{\raggedright\arraybackslash}p{0.24\textwidth}
>{\raggedright\arraybackslash}p{0.24\textwidth}
>{\raggedright\arraybackslash}p{0.20\textwidth}
>{\raggedright\arraybackslash}p{0.22\textwidth}@{}}
\caption{Named omissions and reviewer risks at the 25 August 2026 snapshot.}
\label{tab:omission-register}\\
\toprule
Source or source family & Why its absence could change the proposal & Current
disposition & Next action \\
\midrule
\endfirsthead
\toprule
Source or source family & Reviewer risk & Current disposition & Next action \\
\midrule
\endhead
Gopen and Swan, \emph{The Science of Scientific Writing} & The earlier draft
discussed readability without the best-known account of reader expectations in
scientific prose. & Included, bounded: conceptual editorial model, not efficacy
evidence. & Independent full-text check and backward/forward citation chase. \\
Oppenheimer, \emph{Consequences of Erudite Vernacular} & Claims about needless
complexity otherwise rested on intuition rather than experimental evidence. &
Included, bounded: adjacent processing-fluency evidence. & Classify design,
extract outcomes, and appraise transfer to expert technical readers. \\
Kirchenbauer et al., \emph{A Watermark for Large Language Models} & A blanket
rejection of detection would ignore generation-time provenance. & Included as
a boundary and alternative, not product-efficacy support. & Appraise the threat
model and test whether watermark verification belongs in a future adapter. \\
Sadasivan et al., \emph{Can AI-Generated Text Be Reliably Detected?} & The
detector section would omit a central challenge to reliable post-hoc
classification. & Included as challenge evidence. & Independently inspect the
theory, attacks, corpora, and claims retained in the final journal version. \\
Schriver, \emph{Dynamics in Document Design} & Document-design research may
change how page layout, navigation, and reader testing enter the product. &
Named candidate; not used for a claim. & Retrieve and screen relevant chapters
against the reader-task question. \\
Kellogg, \emph{Training Writing Skills: A Cognitive Developmental Perspective}
& Cognitive accounts may change the burden model for expert revision. & Named
candidate; not used for a claim. & Retrieve, screen, and decide whether it adds
to Flower--Hayes and Bereiter--Scardamalia. \\
Sweller, \emph{Cognitive Load During Problem Solving} & ``Reader burden'' may
otherwise borrow cognitive-load language without an appropriate bridge. &
Named candidate; no cognitive-load claim made. & Search the applied technical-
reading literature before importing the construct. \\
Controlled plain-language studies in technical and public-sector documents &
They may offer direct reader outcomes closer to proposals than student essays.
& Source family undercovered. & Run a dedicated query and retain null or
backfire results as challenge records. \\
Expert peer-review and revision-response studies & They may define acceptance,
rounds, and reviewer disagreement more directly than automated-feedback work. &
Source family undercovered. & Search writing studies, technical communication,
and HCI databases with independent screening. \\
Accessibility and screen-reader studies of mathematical documents & Visual
clarity alone may exclude readers using assistive technology. & Source family
undercovered; no accessibility efficacy claim made. & Add accessibility
fixtures and recruit an accessibility reviewer before deployment. \\
\bottomrule
\end{longtable}

Citation chasing is also a record, not a sentence saying that it occurred. The
five seed groups below are the first load-bearing routes. No row is marked
complete because an independent reviewer has not yet recorded the dispositions.

\begin{table}[H]
\centering
\small
\begin{tabularx}{\textwidth}{@{}p{0.25\textwidth}p{0.22\textwidth}p{0.22\textwidth}L@{}}
\toprule
Seed & Backward chase & Forward chase & Completion rule \\
\midrule
Gopen--Swan & Pending & Pending & Screen foundational antecedents and later
tests of reader-expectation principles. \\
Oppenheimer & Pending & Pending & Screen replications, boundary conditions,
and technical-reader applications. \\
Kirchenbauer and related watermarking & Pending & Pending & Record attacks,
robustness studies, and provenance standards. \\
Sadasivan and detector challenges & Pending & Pending & Record rebuttals,
subsequent detectors, and threat-model differences. \\
Meng and Abudalfa systematic reviews & Partial reference-list inspection;
independent dispositions pending & Pending & Reconcile included-study lists
with the Humanvoice candidate and exclusion ledgers. \\
\bottomrule
\end{tabularx}
\caption{Backward and forward snowball ledger. A route is complete only when
retrieved candidates have screening dispositions and provenance.}
\label{tab:snowball-ledger}
\end{table}

\begin{table}[H]
\centering
\small
\begin{tabularx}{\textwidth}{@{}p{0.27\textwidth}p{0.28\textwidth}L@{}}
\toprule
Evidence gate & Current state & What completion changes \\
\midrule
Primary source identity & DOI-bearing bibliography and load-bearing evidence
rows resolve against Crossref; sentence-level support still needs review. &
Bibliographic claims are reproducible without treating identity as support. \\
Full-text support anchors & Partial; several design claims need page or
section anchors. & A reviewer can inspect the exact support rather than trust a
secondary summary. \\
Independent screening & Outstanding. & Selection and omission risk become
measurable. \\
Specialist-domain import & Pass for a bounded RePEc/IDEAS slice; it is not an
exhaustive licensed search. & Economics-domain coverage becomes less locally
optimized, while the slice still requires screening. \\
Design-specific appraisal & Outstanding for load-bearing studies. & Limits,
task fit, and heterogeneity are visible beside each result. \\
\bottomrule
\end{tabularx}
\caption{Evidence work still required before a strong external literature claim.}
\label{tab:evidence-gates-expanded}
\end{table}

The matrix therefore supports a modest conclusion: existing research justifies
testing a bounded, human-controlled revision workflow and warns against
detector-led or fully automatic rewriting. It does not justify saying that
Humanvoice is already validated, nor that any listed package is the best
choice without a held-out benchmark.

\section{Software benchmark protocol}
\label{app:software-benchmark}

\subsection{The candidate inventory is a set of hypotheses}

The software review names parsers, linters, diff tools, citation utilities,
local inference runtimes, and evaluation frameworks. A candidate enters the
inventory because it might satisfy a contract, not because its repository is
popular or its README uses the right language. The benchmark turns that
hypothesis into an observation.

For every candidate, record version or commit, license, supported formats,
runtime dependencies, maintenance signal, network behavior, model or corpus
provenance, and the exact interface used. A package that passes a fixture but
requires an unapproved remote service is not an operational pass. A package
with an elegant API but no maintained release is a pilot risk. These are
different failure modes and should not be collapsed.

\begin{table}[H]
\centering
\small
\begin{tabularx}{\textwidth}{@{}p{0.21\textwidth}p{0.25\textwidth}p{0.24\textwidth}L@{}}
\toprule
Capability & Candidate family & Fixture question & Default disposition \\
\midrule
Source representation & Pandoc or format-native AST & Are claims, equations,
tables, and citations addressable? & Adapt behind a stable internal model. \\
Build and cross-reference & TeX engine plus log parsers & Can a run expose
undefined references and stable page anchors? & Adopt the existing build;
wrap diagnostics. \\
Protected comparison & latexdiff or structured diff & Are mathematical and
numeric changes distinguishable from prose movement? & Pilot, then extend with
object-aware checks. \\
Style diagnostics & Vale, textlint, proselint, deterministic pattern tools &
Can a finding point to evidence and abstain outside its rule? & Adopt narrow
rules; reject global scores as acceptance. \\
Citations & BibTeX/BibLaTeX tooling and metadata services & Do keys resolve and
does identity remain stable across edits? & Adopt local resolution; remote
lookup only with permission. \\
Model adapter & local runtime or approved provider & Does it preserve the
contract, disclose uncertainty, and obey the privacy boundary? & Pilot behind
an explicit adapter; default to human review. \\
Evaluation & pairwise and rubric frameworks & Can results be replayed with
held-out documents and human labels? & Adapt the measurement layer to the
declared estimand. \\
\bottomrule
\end{tabularx}
\caption{Software capability families and their benchmark questions.}
\label{tab:software-capabilities}
\end{table}

\subsection{Fixture design}

The fixture set should be small enough to run on every commit and rich enough
to expose silent corruption. One fixture contains a numbered equation whose
denominator changes in a child revision. Another contains a citation key that
is renamed but still resolves. A third contains a table whose caption moves
across a page break. A fourth contains a quotation and a code block with
characters that a prose normalizer might strip. The final fixtures are
malformed on purpose: an unclosed environment, a missing bibliography entry,
and an unsupported macro.

For fixture $i$, define the expected protected set $O_i$ and observed set
$\widehat O_i$. The basic preservation rate is
\[
  P_O=\frac{\sum_i |O_i\cap\widehat O_i|}{\sum_i |O_i|},
\]
but report the misses by object kind. A 98 percent aggregate can hide every
equation failure if prose objects dominate the corpus. Also report the
abstention rate $A$ and the false-success rate $F$, where a false success is a
run that reports ``safe'' despite a known protected-object mismatch. The
release rule should treat $F$ as more serious than a visible abstention.

\begin{cardexample}{A benchmark result that changes the architecture}
Suppose a structured parser recovers 100 prose spans and 20 citations but
cannot map displayed equations when a custom macro is present. A text diff
can still produce an attractive report, yet the system has no basis for
claiming protected technical revision. The correct result is an abstention for
those documents and an adapter work package, not a green aggregate score.
\end{cardexample}

\subsection{Adopt, adapt, pilot, skip}

The benchmark disposition has four verbs. \emph{Adopt} means the package's
interface and behavior satisfy a contract under the local privacy and
maintenance conditions. \emph{Adapt} means the underlying capability is
useful but must be wrapped to expose stable object ids or abstention.
\emph{Pilot} means the package is promising but evidence is too thin for a
release dependency. \emph{Skip} means the package's behavior or license is
incompatible with the proposed boundary.

The disposition is revisited when a version, model, or input format changes.
It is not a permanent endorsement. A package can be adopted for deterministic
cross-reference checks and skipped for authorship detection; capabilities, not
repository names, are the unit of judgment.

\subsection{The pacing and ordering adapter}
\label{app:pacing-adapter}

The first implementation can remain small. A standard-library adapter reads
LaTeX, removes comments and drawing environments, preserves paragraph and line
locations, and emits a JSON profile. It runs the same heuristic over a target
and the exemplars named in the reading brief. The output includes words per
concept, singleton share, repeated occurrences, burst paragraphs, and a
trailing live-load window. It also emits a first-use list for terms that lack
an expansion, appositive gloss, or taught-later pointer.

The adapter accepts two explicit configuration lists. The first is vocabulary
the reader already owns; the second contains exemption ranges for previews,
maps, decision boxes, and reference tables. A term in either list is not
silently deleted from the profile. The output records the reason it was not
raised as a first-use question. This makes a later change in the reading brief
replayable.

The adapter's result is a worklist, not a quality score. A typical finding
contains a file, line, paragraph excerpt, observed metric, exemplar comparison,
and one question such as ``Should this unit take on fewer named things?'' or
``Can the ordinary object appear before its acronym?'' The author can keep,
repair, or dismiss it. A dismissal remains in the run record so precision can
be estimated without pretending that every signal was correct.

The comparison layer has one further safeguard. When a prose revision adds a
sentence containing a number, a superlative, a comparison, or a causal verb,
the adapter asks for either a citation, a pointer to an existing protected
claim, an elementary calculation, or an author explanation. The rule catches
the near-miss in which a library gloss acquired the unsupported phrase ``most
widely adopted.'' It does not ban ordinary prose and it does not decide that
the added sentence is false. It makes the new assertion visible while the
author still remembers why it was written.

\begin{cardexample}{A safe prose-diff question}
Parent: ``The package provides a sampler.''\\
Child: ``The package provides the most widely adopted sampler.''\\
\textbf{Observed:} a superlative was added without a citation or protected
claim pointer.\\
\textbf{Question:} which source supports ``most widely adopted,'' or should
the sentence retain the narrower description?
\end{cardexample}

The adapter is useful only while its burden is measured. Retain a rule family
when the median minutes saved over the incumbent inspection are positive on
held-out units, subject to the protected-object veto. Otherwise keep the
fixture and retire the rule. This is a local operating decision, not evidence
that the detector measures comprehension.

\section{Reader instruments and scoring forms}
\label{app:reader-instruments}

\subsection{The reconstruction task without project history}

The primary reader instrument should fit the actual decision. Give the reader
the rendered document, the role statement, and the decision question. Do not
give the private source history or the author's intended interpretation. Ask
the reader to answer in their own words:

\begin{enumerate}
\item What problem is the document addressing?
\item What is the proposed mechanism or product?
\item Which observation, calculation, or source most supports the proposal?
\item What remains uncertain or outside the evidence?
\item What action is being requested, from whom, and under what conditions?
\end{enumerate}

After the answers, ask the reader to mark the first place where they had to
stop and reconstruct a missing relation. Record the page or section, quote no
more than two sentences, and describe what they expected to find there. A
reader may report ``I kept reading but could not tell why the algorithm was
named''; that observation is useful even when the five answers are correct.
The instrument also records whether the stopping point was in the main route,
an appendix, a figure, or a table. These distinctions tell the author whether
the remedy is an explanation, a map, or a navigation aid.

Score each answer for presence, correctness, and location of support. A
reader who gives the right action for the wrong reason is not equivalent to a
reader who reconstructs the argument. Ask for confidence and elapsed time,
then ask for one sentence explaining what would change the decision. The last
question distinguishes a reader who has understood the uncertainty from one
who has merely recognized the requested answer.

\subsection{Rubric and adjudication}

Use a four-point rubric to keep anchors visible:

\begin{table}[H]
\centering
\small
\begin{tabularx}{\textwidth}{@{}p{0.22\textwidth}p{0.17\textwidth}L@{}}
\toprule
Dimension & Score & Anchor \\
\midrule
Problem reconstruction & 0--3 & 0 absent; 1 vague topic; 2 correct problem
with missing boundary; 3 correct problem and population. \\
Mechanism reconstruction & 0--3 & 0 absent; 1 list of features; 2 plausible
sequence with a gap; 3 sequence tied to the reader's decision. \\
Evidence and scope & 0--3 & 0 unsupported; 1 citation or number only; 2
support found but limit missed; 3 support and limit both stated. \\
Action and uncertainty & 0--3 & 0 no action; 1 action guessed; 2 action and
some condition; 3 action, owner, condition, and change trigger. \\
\bottomrule
\end{tabularx}
\caption{Reader reconstruction rubric for the feasibility case.}
\label{tab:reader-instrument}
\end{table}

Two adjudicators independently score a sample and record the sentence or page
that supports each score. Disagreements are resolved by discussion only after
the independent scores are stored. A model judge may suggest a score, but it
cannot replace the human anchor. Report agreement by dimension and preserve
the difficult cases; those cases are often the most informative design tests.

\subsection{Author burden and revision log}

The author instrument records each finding opened, the time spent, the action
taken, and whether a later reader request traced back to that finding. It also
records work that the product caused but that would not have occurred in the
incumbent workflow, such as rechecking a protected equation after an
unnecessary warning. This prevents a reduction in feedback rounds from being
celebrated when total review time has increased.

\begin{longtable}{@{}p{0.22\textwidth}p{0.30\textwidth}p{0.34\textwidth}@{}}
\caption{Minimum event log for measuring workflow burden.}
\label{tab:author-log}\\
\toprule
Event & What to record & Interpretation \\
\midrule
\endfirsthead
\toprule
Event & What to record & Interpretation \\
\midrule
\endhead
Finding opened & id, category, location, timestamp & Exposure to the queue,
not acceptance. \\
Finding accepted & action, rationale, child run & A revision decision, not a
quality label. \\
Finding rejected & rationale and confidence & Possible diagnostic false
positive or legitimate author choice. \\
Abstention & object, reason, fallback & Evidence about the product boundary. \\
Rollback & parent/child runs and reason & Safety or usability failure to
investigate. \\
Reader request & page or section, quoted span, missing relation, requested
change, substantive status & Outcome that links author work to reader work and
preserves the reader's stopping point. \\
\bottomrule
\end{longtable}

\subsection{Disclosure and consent}

The reader should know what assistance shaped the document when that fact can
affect trust or interpretation. The disclosure form should state whether a
model generated questions, whether a human accepted every revision, whether
remote processing occurred, and which source or data were retained. The form
should not imply that model involvement makes a document unreliable; it should
give the reader enough information to interpret the provenance and ask for a
different review path.

For a feasibility study, consent should cover the document's sensitivity,
recorded response time, quotations used in internal reports, and the right to
withdraw a response before analysis. A reader may decline a model-assisted
condition without forfeiting the ordinary review. This is both an ethical
requirement and a useful continuity test.

\section{Cost model and sensitivity worksheets}
\label{app:cost-worksheets}

\subsection{Units and inputs}

The economic case is about review work avoided or improved, not about tokens
or pages alone. The worksheet should use observed units:

\begin{table}[H]
\centering
\small
\begin{tabularx}{\textwidth}{@{}p{0.15\textwidth}p{0.24\textwidth}L@{}}
\toprule
Symbol & Input & Source for the first estimate \\
\midrule
$N$ & Documents or document units reviewed per period & Repository history or
prospective intake count. \\
$r_0$ & Incumbent review rounds per unit & Version and meeting records. \\
$r_1$ & Review rounds with Humanvoice & Feasibility or comparative run. \\
$t_a$ & Author minutes per round & Time log, not a guessed hourly rate. \\
$t_r$ & Reader minutes per round & Timed reader instrument. \\
$w_a,w_r$ & Loaded author and reader rates & Sponsor's accounting convention. \\
$p$ & Probability a unit reaches the decision stage & Intake and completion
records. \\
$c$ & Cost of a protected-object incident or missed claim & Scenario range
agreed before observing the result. \\
$K$ & MVP and recurring operating cost & Work-package estimate and measured
runtime. \\
\bottomrule
\end{tabularx}
\caption{Observable inputs for the economic worksheet.}
\label{tab:cost-inputs}
\end{table}

The value of a workflow over a period can be written as
\[
  V=Np\{(r_0-r_1)(w_at_a+w_rt_r)-\Delta c\}-K,
\]
where $\Delta c$ is the change in expected incident cost. The expression is
deliberately transparent. It does not assume that every avoided round is a
benefit: if a lower round count comes from skipping necessary review, the
incident term should rise. It also permits a negative result when the product
adds author burden or when volume is too small to amortize the build.

\subsection{Worked sensitivity}

Consider a small portfolio with $N=24$ units per year and $p=0.75$. Suppose
the incumbent averages four rounds and the measured workflow averages three,
with 35 author minutes and 25 reader minutes per round. At loaded rates of
\$150 and \$250 per hour, the gross avoided labor is approximately
\[
  24(0.75)(1)\left(150\frac{35}{60}+250\frac{25}{60}\right)
  =\$2{,}025.
\]
This is an illustration, not a forecast. If the MVP costs \$30,000, the first
year cannot be justified by this small portfolio on labor savings alone. The
case changes with a larger volume, a higher cost of a delayed decision, or a
measured reduction in protected-object incidents. The sponsor should see this
arithmetic before hearing a payback period.

\begin{table}[H]
\centering
\small
\begin{tabularx}{\textwidth}{@{}p{0.25\textwidth}rrrL@{}}
\toprule
Scenario & $N$ & $r_0-r_1$ & $K$ & Interpretation \\
\midrule
Small portfolio & 24 & 1.0 & \$30,000 & Learning case; labor alone does not
repay the build. \\
Institutional volume & 180 & 0.8 & \$30,000 & Break-even depends on actual
rates, reader time, and incident avoidance. \\
High-consequence review & 60 & 0.5 & \$30,000 & A modest round reduction can
be valuable if the incident-cost range is credible. \\
Null workflow & 60 & 0.0 & \$30,000 & Stop or reposition unless another
benefit is observed. \\
\bottomrule
\end{tabularx}
\caption{Illustrative sensitivity cases; values are placeholders for observed inputs.}
\label{tab:cost-sensitivity-expanded}
\end{table}

The worksheet should vary one input at a time and then show a joint low,
central, and high case. Do not select the central case after seeing which
scenario crosses break-even. Pre-register the ranges, show the result as a
range, and state which observation would move the project from learning to
deployment. A financier can then challenge the assumptions directly rather
than debate an opaque return percentage.

\subsection{What the model cannot price}

Some benefits and harms are difficult to monetize: a reader's trust, a
reputational error, a lost technical qualification, or the option value of a
reusable source map. They should not disappear from the proposal, but they
should be reported as non-financial outcomes with an owner and an observation
plan. A qualitative benefit is not a license to inflate the financial case.

\section{Reproducible build and source map}
\label{app:reproducibility}

\subsection{What a reader needs}

The reader-facing proposal should stand on its own. The implementation team
needs more: source paths, build commands, dependency versions, bibliography
identity, and a way to compare a regenerated PDF with the recorded one. Keep
those functions separate. The main text explains why reproducibility matters;
the source map tells an operator how to exercise it.

\begin{table}[H]
\centering
\small
\begin{tabularx}{\textwidth}{@{}p{0.24\textwidth}p{0.31\textwidth}L@{}}
\toprule
Artifact & Reproducibility field & Why it is retained \\
\midrule
Canonical source & path, content hash, route inputs, date & Identifies the
reader-facing argument. \\
Build & engine, flags, source-date setting, bibliography pass & Recreates
pagination and references as far as the environment permits. \\
Evidence source & key, URL or DOI, inspected version, anchor, claim class &
Lets a reviewer inspect support and limits. \\
Fixture & source text, expected objects, expected failure, license & Makes a
software result replayable and legally usable. \\
Reader run & document hash, contract, rubric, response, disclosure & Connects
the outcome to the exact document read. \\
Decision record & threshold, observed result, owner, next action & Shows why
the project proceeded, changed, or stopped. \\
\bottomrule
\end{tabularx}
\caption{Source map for a reproducible proposal and feasibility run.}
\label{tab:source-map}
\end{table}

\subsection{Build checks and interpretation}

A clean LaTeX build is necessary but weak evidence. The reproducibility check
should run the structural diagnostics, compile the source twice, inspect the
log for unresolved references and severe box problems, extract text for a
word-count sanity check, and render representative pages for visual review.
The PDF hash is useful only when the environment and source-date setting are
recorded; a changed TeX engine can produce a different hash without a changed
argument.

The visual sample should include the title and decision page, a dense table, a
long equation, a repair pair, a part transition, and an appendix longtable.
Inspect both the page before and after each part break. A chapter that ends
with a heading and two lines of text may be technically valid but rhetorically
broken. Likewise, a table that fits after shrinking the font may be unreadable
to the intended executive reader.

\subsection{Status boundary}

This appendix records a reproducible author-repaired draft. It does not turn
the repository's automated status into a human endorsement. The current
evidence gates remain release-blocked until independent screening,
source-level verification, held-out software benchmarks, and external reader
review are complete. That statement belongs here because it protects the
interpretation of the proposal; a procedural history of how an agent produced
the files does not belong in the reader's opening.

\begin{plainbox}{The final question for an implementer}
Can a person who did not write this proposal take one source snapshot, run the
declared commands, open a finding, inspect the protected object, record an
author decision, give the resulting packet to a named reader, and recover the
evidence behind the next investment decision? If the answer is no, the next
work is not more prose. It is the missing record, fixture, or interface.
\end{plainbox}
